% !PRINTED 301
\weberpage{301}

But $\gamma$ can be further determined so that

\weq{12}{A^{\gamma-\lambda}CA^{-\gamma}=C,}

and for this one obtains the condition

\[ \gamma(g-g^{-1})\equiv\lambda g\quad(\mathrm{mod}\ p), \]

which, since $p>3$, can always be satisfied; but then, on account of (11),

\[ C=A^{\gamma}SA^{-\gamma}. \]

Thus $C$ occurs in the divisor $A^{\gamma}\mathfrak{K} A^{-\gamma}$ conjugate to $\mathfrak{K}$, and we may therefore, by putting this group in place of $\mathfrak{K}$, assume without loss of generality that $\mathfrak{K}$ itself contains the substitution $C$.

Now the substitution

\[ B=\begin{pmatrix}0, & 1\\ -1, & 0\end{pmatrix} \]

may either occur in $\mathfrak{K}$ or not occur. In the first case $\mathfrak{K}$ also contains all the substitutions composed of $B$ and $C$, which are all of one of the two forms:

\weq{13}{\begin{pmatrix}g^{r}, & 0\\ 0, & g^{-r}\end{pmatrix},\qquad r=0,1,\dots\frac{p-3}{2},}

\weq{14}{\begin{pmatrix}0, & g^{-r}\\ -g^{r}, & 0\end{pmatrix},}

and whose number amounts to $p-1$. In the first case, then, $\mathfrak{K}$ contains all substitutions of the form (13), (14), in the second all substitutions of the form (13) and \gesperrt{none} of the form (14).

Now let

\weq{15}{V=\begin{pmatrix}\alpha, & \beta\\ \gamma, & \delta\end{pmatrix}}

be a substitution contained in $\mathfrak{K}$ which is contained neither in the form (13) nor in the form (14), and for which therefore neither $\alpha$ and $\delta$, nor $\beta$ and $\gamma$, are simultaneously congruent to $0$. Then all the substitutions of the form

\weq{16}{C^{r}VC^{s}=\begin{pmatrix}\alpha g^{r+s}, & \beta g^{r-s}\\ \gamma g^{-r+s}, & \delta g^{-r-s}\end{pmatrix},}

when $r,s$ both run through the sequence of numbers:

\[ 0,1,2,\dots\frac{p-3}{2} \]

are contained in $\mathfrak{K}$ and are all \gesperrt{different from one another}. For if two among the substitutions (16) are equal to one another, then we must have
%CONT>>
